\section{Criteri per la scelta dei casi di test}

Collegandoci alle tecniche descritte nella Sezione~2, in questa sezione definiamo i criteri con cui selezioniamo i dati e i casi di prova. Non è pensabile testare l'app in modo esaustivo: le combinazioni tra stati della rete, permessi, ruoli e sensori (GPS, fotocamera, bussola) sono troppe. Per questo motivo abbiamo stabilito dei criteri precisi per individuare i test più significativi.

In linea con le definizioni del corso, distinguiamo chiaramente tra:
\begin{itemize}
    \item \textbf{Dati di test:} l'insieme dei valori di input forniti per testare il sistema;
    \item \textbf{Casi di test:} l'insieme degli input unito agli output attesi corrispondenti, formalizzati come coppie:
    \[
    TC = \langle \text{input}, \: \text{output atteso} \rangle
    \]
\end{itemize}
Nella pratica del progetto, per collegare ogni test ai requisiti del D3, corrediamo ciascuna coppia con un identificativo univoco e il requisito tracciato.

L'obiettivo è scegliere un criterio $C$ che risulti sia \textbf{affidabile} (i test scelti da $C$ danno sempre lo stesso risultato a parità di condizioni) sia \textbf{valido} (se c'è un errore nel programma, esiste almeno un test scelto da $C$ in grado di individuarlo). Come ricordano il teorema di Goodenough-Gerhart e Dijkstra (\emph{«il testing può mostrare la presenza di errori, ma mai la loro assenza»}), nessun criterio garantisce la correttezza assoluta su programmi non banali; perciò combiniamo approcci funzionali, strutturali e di verifica statica.

\subsection{Criteri funzionali (Black-Box)}
I test funzionali derivano direttamente dai requisiti del D3, senza guardare il codice sorgente:

\begin{itemize}
    \item \textbf{Classi di equivalenza:} per ogni dato in ingresso partizioniamo i possibili valori in classi valide e non valide. Scegliamo poi un test per ciascuna classe valida e un test per ciascuna classe non valida presa singolarmente, evitando casi ridondanti.
    \item \textbf{Analisi dei valori limite (BVA):} per i valori numerici o continui (come le coordinate in \texttt{RF-16}, la pendenza della giacitura in \texttt{RF-19} o i parametri chimici in \texttt{RF-11}), oltre ai valori tipici testiamo i valori di frontiera dell'intervallo $[min, max]$:
    \[
    \{min-1, \: min, \: min+1, \: max-1, \: max, \: max+1\}
    \]
    \item \textbf{Pairwise per parametri combinati:} per le funzioni che dipendono da più variabili (\texttt{RF-04} e \texttt{RF-10}), usiamo l'approccio pairwise a due vie. Invece di provare tutte le 24 combinazioni tra Ruolo, Rete e Disciplina, selezioniamo un sottoinsieme di test in cui ogni coppia di valori compare almeno una volta.
    \item \textbf{Transizione di stato:} per la gestione offline (\texttt{RF-08}, \texttt{RF-09}, \texttt{RF-10}) verifichiamo i cambi di stato del dato: \emph{Bozza $\to$ Salvataggio locale $\to$ Coda di invio $\to$ Sincronizzato}, testando anche cosa succede se il processo si interrompe a metà.
\end{itemize}

\subsection{Criteri strutturali (White-Box)}
Per il codice Java controllato con JUnit e JaCoCo (Sezione~1), valutiamo la bontà della suite di test in base alla copertura del codice:

\begin{itemize}
    \item \textbf{Copertura delle decisioni (Branch Coverage):} puntiamo ad almeno l'\textbf{80\%} sui moduli con più logica: validazione dati (\texttt{RF-11}), calcolo permessi (\texttt{RF-03}, \texttt{RF-04}) e versionamento (\texttt{RF-12}). Ogni ramo di un \texttt{if}, \texttt{switch} o ciclo deve essere percorso sia come vero che come falso.
    \item \textbf{Copertura delle istruzioni (Statement Coverage):} manteniamo una soglia minima dell'\textbf{85\%} per le classi principali e i DAO del database Room.
    \item \textbf{Copertura delle condizioni composte:} quando una regola combina più controlli (ad esempio ruolo, gruppo e connessione in \texttt{RNF-11}), ogni singola condizione booleana deve essere verificata sia vera che falsa.
    \item \textbf{Copertura dei cammini:} limitata alla risoluzione dei conflitti durante la sincronizzazione (\texttt{RF-10}), dove percorrere tutti i rami logici è indispensabile per essere certi di non perdere o sovrascrivere dati per sbaglio.
\end{itemize}

\subsection{Criteri per requisiti non funzionali e verifica statica}
\begin{itemize}
    \item \textbf{Requisiti non funzionali (RNF):} selezioniamo casi di prova per situazioni critiche sul campo: arresto forzato dell'app o spegnimento del telefono mentre si scrive su DB (\texttt{RNF-07}), perdita improvvisa del segnale GPS (\texttt{RNF-08}), revoca dei permessi da impostazioni Android (\texttt{RNF-11}) e rotazione dello schermo mentre si compila una scheda (\texttt{RNF-10}).
    \item \textbf{Criteri di accettazione per l'Inspection:} prima di integrare una pull request nel ramo principale, il codice deve superare la lettura a 4 ruoli (Sezione~1). Non accettiamo codice con risorse rimaste aperte (sensori o cursori DB), eccezioni ignorate o segnalazioni gravi di Android Lint (\texttt{RNF-14}, \texttt{RNF-15}).
\end{itemize}

\subsection{Priorità dei test}
Suddividiamo i test in tre fasce per decidere l'ordine di esecuzione:
\begin{itemize}
    \item \textbf{P1 - Bloccante (Alta):} funzioni indispensabili per lavorare sul campo senza connessione e salvare i dati (\texttt{RF-01}, \texttt{RF-06}, \texttt{RF-08}, \texttt{RF-09}, \texttt{RF-10}, \texttt{RF-11}, \texttt{RF-28}). Per rilasciare il codice, il 100\% di questi test deve passare con zero bug aperti.
    \item \textbf{P2 - Importante (Media):} rilievi con sensori, mappe e filtri (\texttt{RF-14}, \texttt{RF-16}, \texttt{RF-18}, \texttt{RF-19}, \texttt{RF-23}, \texttt{RF-29}). Richiediamo almeno il 90\% di esito positivo; eventuali difetti residui devono avere una soluzione temporanea documentata.
    \item \textbf{P3 - Minore (Bassa):} funzioni secondarie come note vocali o dettagli grafici (\texttt{RF-22}, \texttt{RF-26}, \texttt{RF-27}, \texttt{RF-32}), verificate in base al tempo rimasto prima della consegna.
\end{itemize}
